
\subsubsection{Experimental Questions}

\textbf{EQ1: Is test accuracy variance statistically non-zero and practically detectable?}

Experiment H-E1 trained N=30 seeds $\times$ 2 architectures $\times$ 2 datasets (120 total runs). For each condition:
\begin{itemize}
\item Measured test accuracy after 10 epochs of deterministic training
\item Computed variance $\sigma ^2$ across 30 seeds
\item Tested $H_0$: $\sigma ^{2}=0$ vs. $H_1$: $\sigma ^2$>0 using chi-squared test ($\alpha =0.05)$
\item Assessed practical detectability: $\sigma ^2 \geq$ 0.3\%
\end{itemize}

\textbf{EQ2: Do different seeds create independent weight configurations?}

Experiment H-M1 initialized 30 models with seeds 0-29 per condition (no training), flattened weight tensors, computed 435 pairwise Euclidean distances, and tested $H_0$: mean\_distance=0 vs. $H_1$: mean\_distance>0.

\textbf{EQ3: Do different initializations lead to different local minima?}

Experiment H-M2 trained 30 models per condition, computed final weight pairwise distances, and measured loss coefficient of variation.

\textbf{EQ4: Is N=30 sufficient for stable variance estimation?}

Experiment H-M3 bootstrap resampled 30 test accuracies (B=1000 resamples), computed variance for each sample, constructed 95\% confidence intervals, and measured relative width.

\subsubsection{Hardware and Software}

All experiments used a single NVIDIA H100 NVL GPU. Implementation used PyTorch 2.0. Training, evaluation, and analysis code followed deterministic protocols.

\subsubsection{Reproducibility Controls}

\begin{itemize}
\item Fixed hyperparameters (lr=0.01, momentum=0.9, epochs=10, batch\_size=64)
\item Sequential seeds 0-29 (no cherry-picking)
\item Standard train/test splits from torchvision
\item PyTorch seed control + cudnn.deterministic + fixed data order
\end{itemize}
